\section*{अध्याय \ref{ch:g10:lewis-and-shape} --- लुइस संरचनाएँ और अणुओं की आकृति}

\begin{solution}{exo:g10:lewis-and-shape:1}
दो \omterm{def:g7:atoms-and-molecules:atom}{परमाणुओं} द्वारा साझा किया गया \omterm{def:g9:inside-the-atom:nucleus}{इलेक्ट्रॉनों} का एक युग्म। \omterm{def:g10:electron-shells:valence}{संयोजकता इलेक्ट्रॉनों}
का ऐसा युग्म जो केवल एक \omterm{def:g7:atoms-and-molecules:atom}{परमाणु} का है, साझा नहीं।
\end{solution}

\begin{solution}{exo:g10:lewis-and-shape:2}
H: 1 (अपने द्विक से एक \omterm{def:g9:inside-the-atom:nucleus}{इलेक्ट्रॉन} कम)। C: 4, N: 3, O: 2 (अष्टक से 4, 3 और 2
\omterm{def:g9:inside-the-atom:nucleus}{इलेक्ट्रॉन} कम)।
\end{solution}

\begin{solution}{exo:g10:lewis-and-shape:3}
\chemfig{H-\charge{0=\:,90=\:,270=\:}{Cl}}: क्लोरीन के पास साझा युग्म और तीन
\omterm{def:g10:lewis-and-shape:lone-pair}{एकाकी युग्म} हैं, 8 \omterm{def:g9:inside-the-atom:nucleus}{इलेक्ट्रॉन}।
\end{solution}

\begin{solution}{exo:g10:lewis-and-shape:4}
$5 + 3 \times 1 = 8$ \omterm{def:g9:inside-the-atom:nucleus}{इलेक्ट्रॉन}, 4 युग्म (तीन आबंध, एक \omterm{def:g10:lewis-and-shape:lone-pair}{एकाकी युग्म})।
\end{solution}

\begin{solution}{exo:g10:lewis-and-shape:5}
\ce{CH4}: चतुष्फलक। \ce{NH3}: पिरामिड। \ce{H2O}: मुड़ा हुआ। \ce{CO2}:
रैखिक।
\end{solution}

\begin{solution}{exo:g10:lewis-and-shape:6}
\chemfig{H-\charge{90=\:,270=\:}{S}-H}: पानी की तरह, सल्फ़र के चारों ओर चार समूह,
जिनमें दो \omterm{def:g10:lewis-and-shape:lone-pair}{एकाकी युग्म}: मुड़ा हुआ।
\end{solution}

\begin{solution}{exo:g10:lewis-and-shape:7}
$2 \times 5 = 10$ \omterm{def:g9:inside-the-atom:nucleus}{इलेक्ट्रॉन}, 5 युग्म। एक एकल आबंध और 4 \omterm{def:g10:lewis-and-shape:lone-pair}{एकाकी युग्म} हर
नाइट्रोजन के पास 6 \omterm{def:g9:inside-the-atom:nucleus}{इलेक्ट्रॉन} छोड़ते हैं; दो \omterm{def:g10:lewis-and-shape:lone-pair}{एकाकी युग्मों} को साझा युग्मों में बदलने
से एक \omterm{def:g10:lewis-and-shape:multiple-bond}{त्रिआबंध} और हर \omterm{def:g7:atoms-and-molecules:atom}{परमाणु} पर एक \omterm{def:g10:lewis-and-shape:lone-pair}{एकाकी युग्म} मिलता है:
\chemfig{\charge{180=\:}{N}~\charge{0=\:}{N}}।
\end{solution}

\begin{solution}{exo:g10:lewis-and-shape:8}
\chemfig{H-C(-[2]H)(-[6]H)-\charge{90=\:,270=\:}{O}-H}: कार्बन के चारों ओर चार आबंध,
चतुष्फलकीय; ऑक्सीजन के चारों ओर दो आबंध और दो \omterm{def:g10:lewis-and-shape:lone-pair}{एकाकी युग्म}, मुड़ा हुआ।
\end{solution}

\begin{solution}{exo:g10:lewis-and-shape:9}
\ce{CO2} में कार्बन के पास दो समूह (दो \omterm{def:g10:lewis-and-shape:multiple-bond}{द्विआबंध}) हैं और कोई \omterm{def:g10:lewis-and-shape:lone-pair}{एकाकी युग्म} नहीं: वे
विपरीत दिशाओं में होते हैं। \ce{H2O} में ऑक्सीजन के पास चार समूह (दो
आबंध, दो \omterm{def:g10:lewis-and-shape:lone-pair}{एकाकी युग्म}) चतुष्फलकीय व्यवस्था में हैं: दोनों आबंध एक कोण
बनाते हैं।
\end{solution}

\begin{solution}{exo:g10:lewis-and-shape:10}
\chemfig{C(-[3]H)(-[5]H)=C(-[1]H)-[7]H}: हर कार्बन के चारों ओर तीन
समूह (एक \omterm{def:g10:lewis-and-shape:multiple-bond}{द्विआबंध}, दो एकल आबंध): समतल त्रिकोण, कोण लगभग
\ang{120}।
\end{solution}

\begin{solution}{exo:g10:lewis-and-shape:11}
केंद्र में कार्बन, क्लोरीन \omterm{def:g7:atoms-and-molecules:atom}{परमाणुओं} से चार एकल आबंध, हर क्लोरीन पर तीन
\omterm{def:g10:lewis-and-shape:lone-pair}{एकाकी युग्म}: कार्बन के चारों ओर चार समूह, एक चतुष्फलक।
\end{solution}

\begin{solution}{exo:g10:lewis-and-shape:12}
एथेनॉल, \chemfig{H-C(-[2]H)(-[6]H)-C(-[2]H)(-[6]H)-\charge{90=\:,270=\:}{O}-H}, और
मेथॉक्सीमेथेन,
\chemfig{H-C(-[2]H)(-[6]H)-\charge{90=\:,270=\:}{O}-C(-[2]H)(-[6]H)-H}।
\end{solution}

\begin{solution}{exo:g10:lewis-and-shape:13}
\omterm{def:g10:lewis-and-shape:lone-pair}{एकाकी युग्म} \omterm{def:g10:lewis-and-shape:lone-pair}{आबंधी युग्म} से थोड़ी ज़्यादा जगह घेरता है और आबंधों को पास-पास धकेल
देता है। अमोनिया में एक \omterm{def:g10:lewis-and-shape:lone-pair}{एकाकी युग्म} है, पानी में दो: मेथेन से अमोनिया और फिर पानी तक
कोण घटता जाता है।
\end{solution}

\begin{solution}{exo:g10:lewis-and-shape:14}
\chemfig{H-C~\charge{0=\:}{N}}: कार्बन के चारों ओर दो समूह, रैखिक।
\chemfig{H-C(-[6]H)=\charge{45=\:,315=\:}{O}}: कार्बन के चारों ओर तीन समूह, समतल
त्रिकोण।
\end{solution}

\begin{solution}{exo:g10:lewis-and-shape:15}
$3 \times 6 = 18$ \omterm{def:g9:inside-the-atom:nucleus}{इलेक्ट्रॉन}, 9 युग्म:
\chemfig{\charge{180=\:,270=\:}{O}=\charge{90=\:}{O}-\charge{0=\:,90=\:,270=\:}{O}}। केंद्रीय ऑक्सीजन
के पास तीन समूह हैं (एक \omterm{def:g10:lewis-and-shape:multiple-bond}{द्विआबंध}, एक एकल आबंध, एक \omterm{def:g10:lewis-and-shape:lone-pair}{एकाकी युग्म}): \omterm{def:g7:atoms-and-molecules:molecule}{अणु}
मुड़ा हुआ है।
\end{solution}

\begin{solution}{pb:g10:lewis-and-shape:1}
\textbf{1.} \chemfig{H-C(-[2]H)(-[6]H)-H}, \chemfig{H-\charge{90=\:}{N}(-[6]H)-H},
\chemfig{H-\charge{90=\:,270=\:}{O}-H}।

\textbf{2.} हर एक में चार समूह।

\textbf{3.} मेथेन में कोई नहीं, अमोनिया में एक, पानी में दो।

\textbf{4.} कार्बन चार आबंध बनाता है; ऑक्सीजन दो (उसके बाकी दो
समूह \omterm{def:g10:lewis-and-shape:lone-pair}{एकाकी युग्म} हैं, जिन्हें किट छड़ों से नहीं दिखाती)।

\textbf{5.} \chemfig{H-C(-[2]H)(-[6]H)-C(-[2]H)(-[6]H)-H}: एक छड़
दोनों कार्बन गेंदों को जोड़ती है, और छह हाइड्रोजन गेंदें चाहिए।

\textbf{6.} चतुष्फलक, पिरामिड, मुड़ा हुआ।

\textbf{7.} \omterm{def:g10:lewis-and-shape:lone-pair}{एकाकी युग्म} \omterm{def:g10:lewis-and-shape:lone-pair}{आबंधी युग्मों} से ज़्यादा जगह घेरते हैं और आबंधों को पास-पास
दबा देते हैं; दो \omterm{def:g10:lewis-and-shape:lone-pair}{एकाकी युग्मों} वाले पानी का कोण सबसे छोटा है।

\textbf{8.} कार्बन गेंद के छेद \ang{109.5} बनाते हैं; पानी का वास्तविक कोण
\ang{104.5} है: ऑक्सीजन गेंद को \ang{104.5} पर अपने दो अलग छेद चाहिए।

\textbf{9.} \ang{180}: कार्बन के चारों ओर के दो समूह विपरीत दिशाओं में होते
हैं, और \omterm{def:g7:atoms-and-molecules:molecule}{अणु} रैखिक है।

\textbf{10.} एक ही किनारे पर न होने वाले दो कोने किसी एक फलक के आमने-सामने के कोने
हैं: उनकी दूरी फलक का विकर्ण है, $\sqrt{a^2 + a^2} = a\sqrt{2}$।

\textbf{11.} केंद्र घन के किसी विकर्ण के आधे रास्ते पर है:
$a\sqrt{3}/2$।

\textbf{12.} H--H: $\sqrt{2} \approx \qty{1.41}{cm}$; C--H:
$\sqrt{3}/2 \approx \qty{0.87}{cm}$।

\textbf{13.} हर हाइड्रोजन घन के एक कोने पर है और कार्बन उसके केंद्र पर: हर
कोना केंद्र से समान दूरी पर है।

\textbf{14.} त्रिभुज C\,H$_1$\,H$_2$ समद्विबाहु है (C\,H$_1$ = C\,H$_2$);
उसके शीर्ष से आधार के मध्यबिंदु तक की रेखा आधार पर
लंब होती है।

\textbf{15.} $\sin \widehat{\text{M\,C\,H}_1} = \dfrac{\text{M\,H}_1}{\text{C\,H}_1}
= \dfrac{a\sqrt{2}/2}{a\sqrt{3}/2} = \sqrt{2/3} \approx 0.816$।

\textbf{16.} अर्धकोण $\approx \ang{54.7}$; कोण H--C--H $\approx
2 \times 54.7 = \ang{109.5}$।

\textbf{17.} यह सारणी वाले चतुष्फलक का \ang{109.5} है।

\textbf{18.} छेद एक-दूसरे से \ang{109.5} पर किए जाने चाहिए।
\end{solution}
